\section{Verification of the mixed-tail identities}
\label{tail:verification-section}

The script \texttt{verification/verify\_mixed\_tails.py} records
finite exact checks and a numerical route independent of the moment
formula. Its numerical tail approximation has an explicit analytic
remainder estimate. Ordinary multiprecision arithmetic is used, so
the recorded floating-point computations are not asserted to be
interval certificates. The identities are established by the
proofs above.

\subsection{A proved truncation envelope}

For real \(x>0\), put
\[
 T_p(x)=\zeta(p,x+1/2),\qquad
 Q_{p,K}(x)=\sum_{j=0}^{K}c_{p,j}x^{1-p-2j},\qquad
 R_{p,K}=\frac{2(p+2K)!}{(p-1)!(2\pi)^{2K+2}}.
\]
Then
\begin{equation}\label{tail:individual-bound}
 0<T_p(x)\le\frac{x^{1-p}}{p-1},\qquad
 |T_p(x)-Q_{p,K}(x)|\le R_{p,K}x^{-p-2K-1}.
\end{equation}
Here and below the bound is needed only to substantiate numerical
verification, not as the research objective.

To prove it, let \(g(u)=u/(2\sinh(u/2))\). The standard
hyperbolic-cosecant partial fractions \cite[Eq.~4.36.5]{DLMF}
give
\[
 g(u)=1+2\sum_{m=1}^{\infty}
             (-1)^m\frac{u^2}{u^2+(2\pi m)^2}.
\]
Expanding each rational summand through degree \(2K\), and using
the alternating-series estimate on the remainder, gives for real
\(u>0\)
\[
 \left|g(u)-\sum_{j=0}^{K}
        \frac{B_{2j}(1/2)}{(2j)!}u^{2j}\right|
 \le\frac{2u^{2K+2}}{(2\pi)^{2K+2}}.
\]
The identity
\[
 T_p(x)=\frac1{(p-1)!}\int_0^{\infty}
             u^{p-2}e^{-xu}g(u)\,\dd u
\]
now proves both parts of \eqref{tail:individual-bound}; the
positivity bound uses \(0<g(u)\le1\).

Fix \(N\ge1\), put \(a=N+1/2\), and choose \(K\ge K_r\).
Define
\[
 M_{p,K}(a)=\sum_{j=0}^{K}|c_{p,j}|a^{-2j},\qquad
 k_{p,q,\ell}^{[K]}
   =\sum_{\substack{0\le j\le K\\0\le\ell-j\le K}}
                         c_{p,j}c_{q,\ell-j}.
\]
The finite computational expression is
\begin{align}\label{tail:finite-evaluation}
 V_{p,q,r}^{N,K}={}&
 \sum_{n=0}^{N-1}\left\{x_n^{2r}A_p(n)A_q(n)
     -\sum_{\ell=0}^{K_r-1}
            k_{p,q,\ell}x_n^{2r+2-W-2\ell}\right\}\notag\\
 &+\sum_{\ell=K_r}^{2K}
       k_{p,q,\ell}^{[K]}\zeta(W-2+2\ell-2r,a).
\end{align}
All the Hurwitz-zeta orders in the last sum are at least two.
Provided \(W+2K-2r>1\), its error satisfies
\begin{align}\label{tail:evaluation-envelope}
 |D_{p,q}(-2r)-V_{p,q,r}^{N,K}|
 \le{}&\left\{\frac{R_{p,K}}{q-1}
                   +M_{p,K}(a)R_{q,K}\right\}
                     \zeta(W+2K-2r,a).
\end{align}
One may take the smaller bound obtained by interchanging \(p,q\).
Indeed, \eqref{tail:individual-bound} gives
\[
 |T_pT_q-Q_{p,K}Q_{q,K}|
 \le\left\{\frac{R_{p,K}}{q-1}
                   +M_{p,K}(a)R_{q,K}\right\}x^{-W-2K}
 \quad(x\ge a).
\]
Multiplication by \(x^{2r}\), summation over the omitted tail,
and expansion of the finite product prove
\eqref{tail:evaluation-envelope}.

\subsection{Recorded checks}

The script uses \(110\) decimal working digits, \(N=64\), and
\(K=30\). It checks \(r=0,1,2,3,5,6\) for each of
\[
 (p,q)=(2,2),(2,4),(4,4),(2,6),(4,6),(6,6).
\]
These include the first regularized values and several higher
regularized values, not only ordinary unsubtracted sums.
Every one of the \(36\) observed discrepancies is below its
analytic truncation envelope. Selected results are
\begin{center}
\begin{tabular}{cccll}
\hline
\(p\)&\(q\)&\(r\)&Absolute discrepancy&Truncation envelope\\
\hline
2&4&2&\(1.4110\,10^{-73}\)&\(1.4471\,10^{-73}\)\\
2&4&3&\(5.9899\,10^{-70}\)&\(6.1431\,10^{-70}\)\\
2&4&6&\(4.6138\,10^{-59}\)&\(4.7314\,10^{-59}\)\\
4&4&6&\(7.2161\,10^{-63}\)&\(7.4014\,10^{-63}\)\\
4&6&6&\(1.8213\,10^{-64}\)&\(1.8709\,10^{-64}\)\\
6&6&6&\(5.1274\,10^{-68}\)&\(5.2690\,10^{-68}\)\\
\hline
\end{tabular}
\end{center}

The exact checks independently expand the finite uncentered tail
products and Faulhaber polynomials and extract their constant
coefficients for six pairs and eleven odd block orders. They also
verify the short \((2,4)\) block formula through \(j=12\) and
recover the previously proved trigamma-square formula through
\(r=12\). The two eliminations in
\eqref{tail:24-eliminations} are checked by exact rational algebra.

Finally, direct quadrature of the generator is compared with its
first \(41\) Taylor terms for the three cases
\[
 (p,q,t)=(2,4,0.7),\quad (2,4,1.3),\quad (4,6,0.8).
\]
Their respective differences are approximately
\(1.1\,10^{-78}\), \(1.2\,10^{-56}\), and
\(2.6\,10^{-76}\), consistent with the omitted convergent Taylor
tail. For both \((2,4)\) checks, direct integral evaluation and
the explicit \(\operatorname{Li}_2\)--\(\operatorname{Li}_6\)
formula agree to better than \(10^{-109}\). Complete numerical
values, software versions, formulas, and truncation envelopes are
recorded in \texttt{verification/mixed\_tail\_results.json}.

